% ---- pg-282 （书页 270） ----

\[
\log(\text{当前用电量})-\log(\text{基准用电量})
\]

对于到上午 9 点、11 点(高峰时间)及下午 1 点、3 点(高峰时间)为止所得到的响应测量值, 产生汇总统计数字如下:

测试组: $n_1=28, \bar{\mathbf{x}}_1=[0.153,-0.231,-0.322,-0.339]'$

控制组: $n_2=58, \bar{\mathbf{x}}_2=[0.151,0.180,0.256,0.257]'$

且
\[
\mathbf{S}_p=\begin{bmatrix}0.804&0.355&0.228&0.232\\0.355&0.722&0.233&0.199\\0.228&0.233&0.592&0.239\\0.232&0.199&0.239&0.479\end{bmatrix}
\]
试完成轮廓分析. 实施用电时间定价方案对电力消费是否有影响? 如果有, 那么, 差别的本质是什么? 试作讨论(对任何统计检验均用显著性水平 $\alpha=0.05$).

\noindent 6.27\quad 作为例 6.14 中爱情与婚姻问题研究的一部分, 对一个由若干名丈夫和妻子组成的样本进行了问卷调查, 请他们回答下列问题:

\noindent (1) 你对伴侣的爱情的“热度”感觉如何?

\noindent (2) 伴侣对你的爱情的“热度”感觉如何?

\noindent (3) 你对伴侣的爱情的“可结伴”水平感觉如何?

\noindent (4) 伴侣对你的爱情的“可结伴”水平感觉如何?

\noindent 回答均采用如下 5 级记分制:

\begin{center}
\begin{tikzpicture}[x=2.2cm,y=1cm]
\foreach \i/\lab in {0/没有,1/很小,2/有些,3/很大,4/非常大}
  \node[above] at (\i,0.05) {\lab};
\draw (0,0) -- (4,0);
\foreach \i in {0,1,2,3,4} \draw (\i,0) -- (\i,-0.18);
\foreach \i/\num in {0/1,1/2,2/3,3/4,4/5}
  \node[below] at (\i,-0.18) {\num};
\end{tikzpicture}
\end{center}

30 名丈夫与 30 名妻子的回答列于\hyperref[tab:6.14]{表 6.14}. 其中:

\noindent $X_1=$ 对问题 1 的 5 级分制回答

\noindent $X_2=$ 对问题 2 的 5 级分制回答

\noindent $X_3=$ 对问题 3 的 5 级分制回答

\noindent $X_4=$ 对问题 4 的 5 级分制回答

\begin{mvtable}
\phantomsection\label{tab:6.14}\mvtabcap{表 6.14\quad 配偶数据}
\begin{tabular}{rrrr|rrrr}
\toprule
\multicolumn{4}{c}{丈夫对妻子} & \multicolumn{4}{c}{妻子对丈夫}\\\midrule
$x_1$ & $x_2$ & $x_3$ & $x_4$ & $x_1$ & $x_2$ & $x_3$ & $x_4$\\\midrule
2 & 3 & 5 & 5 & 4 & 4 & 5 & 5\\
5 & 5 & 4 & 4 & 4 & 5 & 5 & 5\\
4 & 5 & 5 & 5 & 4 & 4 & 5 & 5\\
4 & 3 & 4 & 4 & 4 & 5 & 5 & 5\\
3 & 3 & 5 & 5 & 4 & 4 & 5 & 5\\
3 & 3 & 4 & 5 & 3 & 3 & 4 & 4\\
3 & 4 & 4 & 4 & 4 & 3 & 5 & 4\\
4 & 4 & 5 & 5 & 3 & 4 & 5 & 5\\
4 & 5 & 5 & 5 & 4 & 4 & 5 & 4\\
4 & 4 & 3 & 3 & 3 & 4 & 4 & 4\\
\bottomrule
\end{tabular}
\end{mvtable}
